\documentclass{article}

\usepackage[letterpaper, margin=1in]{geometry}
\usepackage{amsmath,amssymb}

% math conventions, copied from main.tex
\newcommand{\Le}{\leqslant}
\newcommand{\Ge}{\geqslant}
\newcommand{\D}{\mathrm{d}}
\newcommand{\KL}{\mathrm{KL}}
\newcommand{\X}{\mathrm{X}}
\newcommand{\ESS}{\mathrm{ESS}}
\newcommand{\DB}{\mathrm{DB}}

\begin{document}

\begin{center}
    {\large\textbf{Q\&A}}
\end{center}

\paragraph{Q1.} What is the motivation of the X functional?

\paragraph{A1.} The starting point is a stronger object than the forward KL, the
\emph{detailed balance divergence}
\begin{equation*}
    \DB(\mu\|\nu) = \int_{\mathbb R^d} \mu(x)\,
    \biggl| \nabla \log \frac{\mu(x)}{\nu(x)} \biggr| \D x .
\end{equation*}
Matching $\nu$ to $\mu$ is the same as flattening the log-ratio $\log(\mu/\nu)$,
and the gradient measures exactly this variation, so $\DB$ vanishes if and only
if $\mu(x)/\nu(x)$ is constant on the support of $\mu$, that is $\nu = \mu$.
Where the forward KL constrains only the mean of $\log(\mu/\nu)$ under the
target, the gradient in $\DB$ forces the whole ratio $\mu(x)/\nu(x)$ to be
constant. It is closely related to the Fisher information, obtained by squaring
the gradient norm $|\nabla \log(\mu/\nu)|$ in $\DB$.

This choice has one fatal drawback. The gradient $\nabla \log(\mu/\nu)$ involves
the score of the pushforward $\nu$, whose evaluation needs the full Jacobian
$J_G$ of the flow --- not merely its determinant, which the flow already
provides --- and is therefore extremely expensive for a deep map.

The remedy is to replace the gradient, an infinitesimal difference at a single
point, by a finite difference between two points, which yields the first X
functional
\begin{equation*}
    \X_\mu(\mu\|\nu) = \iint_{\mathbb R^d\times\mathbb R^d} \mu(x)\,\mu(y)\,
    \biggl| \log\frac{\mu(x)}{\nu(x)} - \log\frac{\mu(y)}{\nu(y)} \biggr|
    \D x \D y .
\end{equation*}
It keeps the property that motivated $\DB$ --- it vanishes exactly when
$\log(\mu/\nu)$ is constant on the support of $\mu$ --- but removes the gradient:
it reads off only the log-ratio $\log(\mu/\nu)$, which every normalizing flow
provides cheaply, so the expensive Jacobian never appears. The pairwise
difference also cancels the unknown normalizer of $\mu$, which enters the
log-ratio as an additive constant and drops out of the difference.

\paragraph{Q2.} Why is the X regularization defined with an absolute value
rather than a square?

\paragraph{A2.} We did consider the squared variant --- the pairwise difference
under a square $|\cdot|^2$ in place of the absolute value $|\cdot|$ in $\X_\mu$
--- and tested both on the 2D benchmarks. The absolute value was consistently
the better of the two, and it almost never caused the numerical singularities
the square is prone to. We suspect the reason is that the absolute value keeps
$\X_\mu$ closest to the forward KL, on the same linear scale in the log-ratio,
so the two can share a single learning rate; the square grows quadratically and
sits on a different scale.

\paragraph{Q3.} Can $\X_\mu$ be used on its own, without the forward KL, as the
loss?

\paragraph{A3.} No. On the 2D benchmarks, $\X_\mu$ alone does not even make the
flow converge. We conjecture why: the forward KL, driven by Jensen's inequality,
shapes the flow globally, whereas $\X_\mu$ only adjusts the log-ratio pointwise
and locally, and a purely local signal cannot drive the flow to converge. The
forward KL is therefore indispensable, and its weight against $\X_\mu$ is a
genuine trade-off between accuracy and cost; we recommend an equal $1{:}1$
weight as the default.

\paragraph{Q4.} Is the X-regularized forward KL always better, and how does one
choose between it and the plain forward KL in practice?

\paragraph{A4.} Not always. In most of the small-batch cases we tested, the
X-regularized losses were better; at large batch, though, the gap between the
forward KL and its X-regularized versions shrinks. The rule of thumb is simple:
when the batch is small, use the X regularization. And if a multi-GPU machine is
at hand, there is no need to commit to one loss --- train several in parallel
within a single stage (forward KL, FAB, $\KL{+}\X{+}\X$, or even
$\mathrm{FAB}{+}\X{+}\X$), compare them, and carry the best one forward.

\paragraph{Q5.} Where does the mode discovery ability of the X functionals come
from, and how does it differ from PT or FAB?

\paragraph{A5.} It comes from the mixture functional $\X_{(\hat\mu+\bar\nu)/2}$,
built on the wide-coverage measure $\hat\mu$ that quench and temper (QT)
prepares before training. PT is a classical Monte Carlo method that leans on
frequent temperature swaps for ergodicity, so it can only move mass between
modes a chain already reaches; FAB, and other methods that merely change the
divergence, inherit the same limit, since their annealing chains must still
cross the barriers and none hunts for a lost mode on its own. QT does find lost
modes proactively, off any chain --- but its samples are coarse, and feeding
them into training directly would bias the target and cost accuracy. The X
functional is what makes them usable: it vanishes whenever the log-ratio is
constant on the support of its weight, whatever the shape of that weight, so
$\hat\mu$ only tells the flow where to place mass, never what density to fit.
The mixture term thus turns the QT data into a clean training signal, pulling
the flow onto the discovered modes while leaving the minimizer $\nu = \mu$
untouched.

\paragraph{Q6.} What is the meaning of the Fisher--Rao gradient flow analysis,
and does it connect to the real flow training?

\paragraph{A6.} Frankly, not much. The Fisher--Rao flow is a continuous-time
idealization on the space of densities, whereas real training takes discrete
Adam steps on finite batches; we claim no faster transient from it. Its real
value, in our view, is the asymptotic accuracy result: under biased annealed
samples, adding $\X_\mu$ provably lowers the accuracy floor at which the bare
forward KL is stuck. And since $\X_\mu$ reuses the forward KL's own samples and
forward pass at negligible cost, it is about as close to a free lunch as one
finds --- a cheap term that simply buys accuracy.

\end{document}
